% Source-bound reconstruction of the corrected/expanded Abramsky--Jung author PDF.
\section{Domain theory and recursive semantic models}
\label{sec:p10-domain-theory-recursive-models}

Recursive values and recursive types pose different fixed-point problems.
The first asks for an element \(d=f(d)\) in a given ordered structure;
the second asks for a structure and an actual isomorphism
\(D\cong F(D)\), possibly with contravariant occurrences of \(D\).
This treatment follows Abramsky and Jung's \emph{Corrected and expanded
version} of \emph{Domain Theory}, with its own page and result numbering
\cite[pp.~1,5--7]{AbramskyJungCorrectedExpandedDomainTheory}.
That author presentation retains numbering while changing and explaining
the published text; it is not identified with the printed 1994 chapter,
the separate PostScript manifestation, or other semantic-domain sources.

\subsection{Directed convergence and approximation}

\begin{hbconvention}{def:p10-domain-order-convention}{Information order and nonempty directed sets}
Write \(x\sqsubseteq y\) for the information order of a poset, not for
an independently given numerical order. A subset \(A\) is
\emph{directed} when it is nonempty and every two members have an upper
bound belonging to \(A\). Consequently every finite subset of \(A\)
has an upper bound in \(A\). Write \(\bigsqcup^{\uparrow}A\) for its
supremum when this exists, and \(\downarrow x=\{y:y\sqsubseteq x\}\).
A \emph{dcpo} is a poset with suprema of all directed subsets. A
\emph{pointed} dcpo has a least element \(\bot\). Pointedness is an
additional condition: a nonempty discrete poset is a dcpo but generally
has no least element. The empty poset is allowed as a dcpo, but cannot
be pointed. All constructions here use set-sized objects and maps;
claims explicitly involving countable unions or enumerations use ordinary
set theory with choice. No reduction of arbitrary directed sets to
countable chains is assumed
\cite[Definitions~2.1.1,2.1.3,2.1.8,2.1.13, pp.~10--15]
{AbramskyJungCorrectedExpandedDomainTheory}.
\end{hbconvention}

\begin{hbdefinition}{def:p10-domain-continuous-map}{Scott-continuous and strict maps}
For dcpo's \(D,E\), a map \(f:D\to E\) is Scott-continuous when it is
monotone and
\[
 f(\bigsqcup^{\uparrow} A)=\bigsqcup^{\uparrow}_{a\in A}f(a)
                    \qquad(A\subseteq D\text{ directed}).
\]
A map between pointed dcpo's is \emph{strict} when
\(f(\bot_D)=\bot_E\). Let \([D\longrightarrow E]\) be all
Scott-continuous maps, ordered pointwise, and
\([D\mathrel{\mathop{\longrightarrow}^{\bot!}}E]\) the strict ones.
The category \(\mathsf{DCPO}\) has all dcpo's and Scott-continuous maps;
\(\mathsf{DCPO}_{\bot}\) restricts objects to pointed dcpo's;
\(\mathsf{DCPO}_{\bot!}\) also restricts maps to strict ones.
These are different categories, not different spellings of one category
\cite[Definition~2.1.17; pp.~15,41]{AbramskyJungCorrectedExpandedDomainTheory}.
\index{dcpo}\index{Scott continuity}\index{strict map!domain theory}
\end{hbdefinition}

\begin{hbdefinition}{def:p10-domain-way-below}{Approximation, bases and domains}
For \(x,y\in D\), put \(x\ll y\) when, for every directed
\(A\subseteq D\),
\[
 y\sqsubseteq\bigsqcup^{\uparrow}A
       \quad\Longrightarrow\quad
       (\exists a\in A)\ x\sqsubseteq a.
\]
Write \(\mathord{\downarrow\!\downarrow}y=\{x:x\ll y\}\),
\(\mathord{\uparrow\!\uparrow}x=\{y:x\ll y\}\), and
\(K(D)=\{x:x\ll x\}\). Members of \(K(D)\) are compact elements.
A subset \(B\subseteq D\) is a basis when, for every \(x\),
\(B_x=B\cap\mathord{\downarrow\!\downarrow}x\) contains a directed
subset with supremum \(x\). This is the source's basis definition;
directedness of the entire \(B_x\) is a consequence, not a second
unmentioned premise. A continuous domain is a dcpo with a basis; an
algebraic domain has a basis of compact elements. The prefixes
\(\omega\)-continuous and \(\omega\)-algebraic mean that a countable
basis exists, respectively that \(K(D)\) is a countable basis.
Here ``domain'' includes approximation as well as convergence
\cite[Definitions~2.2.1,2.2.3,2.2.6, pp.~18--19]
{AbramskyJungCorrectedExpandedDomainTheory}.
\index{way-below relation}\index{continuous domain}\index{algebraic domain}
\end{hbdefinition}

\begin{hbtheorem}{thm:p10-domain-basis}{proved=annex:p10.domain.basis}{Basis properties and interpolation}
In any dcpo, \(x\ll y\) implies \(x\sqsubseteq y\), and
\(x'\sqsubseteq x\ll y\sqsubseteq y'\) implies \(x'\ll y'\).
For a basis \(B\), every \(B_x\) is directed with supremum \(x\),
and \(K(D)\subseteq B\). Hence continuity is equivalent to
\(x=\bigsqcup^{\uparrow}\mathord{\downarrow\!\downarrow}x\) for
every \(x\); algebraicity is equivalent to the same assertion with
\(K(D)\cap\downarrow x\).
If \(D\) is continuous, \(M\subseteq D\) is finite, and every
\(m\in M\) satisfies \(m\ll y\), there is \(b\in B\) with
\(m\ll b\ll y\) for all \(m\in M\), including the empty \(M\).
In particular,
\[
 \mathord{\downarrow\!\downarrow}\bigsqcup^{\uparrow}A
       =\bigcup_{a\in A}\mathord{\downarrow\!\downarrow}a
\]
for directed \(A\) in a continuous domain
\cite[Propositions~2.2.2,2.2.4,2.2.7; Lemma~2.2.15;
Corollary~2.2.16, pp.~18--25]{AbramskyJungCorrectedExpandedDomainTheory}.
\end{hbtheorem}

\begin{hbtheorem}{thm:p10-domain-basic-examples}{proved=annex:p10.domain.examples}{Finite information and genuinely directed examples}
The powerset \(\mathcal P(S)\), ordered by inclusion, is algebraic;
directed suprema are unions, compact elements are exactly finite subsets,
and \(A\ll B\) exactly when \(A\) is finite and \(A\subseteq B\).
For uncountable \(S\), its directed family of finite subsets is not a
chain and has no cofinal chain.
The partial functions \([X\rightharpoonup Y]\), ordered by graph
inclusion, form a pointed algebraic dcpo. Their directed suprema are
graph unions and their compact elements are the finite partial functions.
The interval \([0,1]\), with its numerical order, is continuous but
not algebraic: \(x\ll y\) exactly when \(x=0\) or \(x<y\), and
\(K([0,1])=\{0\}\). A countable basis is the rational points in
\([0,1]\). A dcpo need not be continuous: two incomparable increasing
\(\omega\)-chains with a common supremum, even after adding bottom,
have only bottom as an approximant to their nonbottom elements.
These reconstruct the source examples and Exercises~2.3.9(6),(19),
with the additional pointed two-chain comparison made explicit
\cite[pp.~13,18--21,31--32]{AbramskyJungCorrectedExpandedDomainTheory}.
\end{hbtheorem}

\begin{hbdefinition}{def:p10-domain-abstract-basis}{Abstract bases and round ideals}
An abstract basis \((B,\prec)\) has a transitive relation such that
for every finite \(M\subseteq B\) and \(x\in B\),
\[
  (\forall m\in M)\ m\prec x
       \quad\Longrightarrow\quad
  (\exists y\in B)\ (\forall m\in M)\ m\prec y\prec x.
\]
An ideal is a nonempty lower subset directed with respect to \(\prec\):
for \(a,b\) in it some \(c\) in it satisfies \(a\prec c\) and
\(b\prec c\). Taking \(a=b\) enforces roundness. Let
\(\operatorname{Idl}(B,\prec)\) be its ideals ordered by inclusion,
and \(i(b)=\{a:a\prec b\}\). Neither reflexivity nor
\(b\in i(b)\) is assumed
\cite[Definitions~2.2.20--21, pp.~26--27]
{AbramskyJungCorrectedExpandedDomainTheory}.
\end{hbdefinition}

\begin{hbtheorem}{thm:p10-domain-ideal-completion}{proved=annex:p10.domain.ideals}{Actual ideal completion and recovery from a basis}
The ideal completion is a continuous dcpo with directed unions as
suprema and basis \(i(B)\). Its approximation relation is
\[
 A\ll A'\quad\Longleftrightarrow\quad
       (\exists b\in A')\ A\subseteq i(b).
\]
For reflexive \(\prec\) it is algebraic. In particular the ideal
completion of a poset \(B\) has compact elements exactly its
principal ideals. If \(B\) is a basis of a continuous domain \(D\),
with \(\prec\) the restricted \(\ll\), the inverse isomorphisms are
\[
 A\longmapsto\bigsqcup^{\uparrow}A,
       \qquad x\longmapsto B_x.
\]
If \(f:D\to E\) is continuous, its values on that basis determine it:
\(f(x)=\bigsqcup^{\uparrow}_{b\in B_x}f(b)\)
\cite[Propositions~2.2.22,2.2.25; Corollary~2.2.26, pp.~27--28]
{AbramskyJungCorrectedExpandedDomainTheory}.
\end{hbtheorem}

\begin{hbdefinition}{def:p10-domain-scott-topology}{Scott-open observations}
A Scott-open subset \(O\subseteq D\) is upper and inaccessible by
directed suprema: \(\bigsqcup^{\uparrow}A\in O\) implies
\(A\cap O\ne\varnothing\). Equivalently its complement is lower
and closed under directed suprema. Write \(\sigma D\) for these
opens. This topology need not be Hausdorff
\cite[Definition~2.3.1, p.~29]{AbramskyJungCorrectedExpandedDomainTheory}.
\end{hbdefinition}

\begin{hbtheorem}{thm:p10-domain-scott-topology}{proved=annex:p10.domain.topology}{Topology, order and continuity agree at the stated scope}
The Scott opens form a topology and recover the order:
\(x\sqsubseteq y\) iff every Scott-open neighborhood of \(x\)
contains \(y\). Thus this topology is \(T_0\).
A map of dcpo's is Scott-continuous exactly when it is continuous
between their Scott topologies. In a continuous domain with basis
\(B\), the sets \(\mathord{\uparrow\!\uparrow}b\), \(b\in B\),
form a basis for \(\sigma D\).
\cite[Propositions~2.3.2,2.3.4--6, pp.~29--31]
{AbramskyJungCorrectedExpandedDomainTheory}.
\end{hbtheorem}

\subsection{Least fixed points as meanings of recursion}

\begin{hbtheorem}{thm:p10-domain-function-space}{proved=annex:p10.domain.functions}{Function dcpo's, evaluation and currying}
For dcpo's \(D,E\), \([D\longrightarrow E]\) is a dcpo whose
directed suprema are pointwise. If \(E\) is pointed it has bottom
the constant \(\bot_E\) map. For pointed \(D,E\), the strict function space is also a
dcpo with pointwise directed suprema.
Finite products are coordinatewise dcpo's, and their approximation
relation is coordinatewise; products of bases are bases.
Evaluation, currying and composition are Scott-continuous, with
\[
 [D\times E\longrightarrow F]\cong[D\longrightarrow[E\longrightarrow F]].
\]
The functorial action on function spaces is
\([f,g](h)=g\circ h\circ f\), contravariant in \(f\)'s object.
This proves cartesian closure of \(\mathsf{DCPO}\), not of all
continuous or algebraic domains
\cite[Proposition~2.1.18; Section~3.2.1--2, pp.~16,42--44]
{AbramskyJungCorrectedExpandedDomainTheory}.
\end{hbtheorem}

\begin{hbtheorem}{thm:p10-domain-least-fixpoint}{proved=annex:p10.domain.fixpoint}{Least fixed point and continuity of its operator}
Let \(D\) be a pointed dcpo and \(f:D\to D\) Scott-continuous.
Then
\[
 \operatorname{fix}(f)=\bigsqcup^{\uparrow}_{n\in\mathbb N}f^n(\bot)
\]
is a fixed point and is below every pre-fixed point
\(x\) with \(f(x)\sqsubseteq x\). The actual assignment
\(\operatorname{fix}:[D\longrightarrow D]\to D\) is Scott-continuous.
For continuous \(F:P\times D\to D\), the parametrized solution
\(p\mapsto\operatorname{fix}(F(p,-))\) is therefore continuous and
satisfies its recursive equation
\cite[Theorem~2.1.19; Exercise~2.3.9(17), pp.~16,32]
{AbramskyJungCorrectedExpandedDomainTheory}.
\end{hbtheorem}

\begin{hbtheorem}{thm:p10-domain-fixpoint-principles}{proved=annex:p10.domain.induction}{Induction and strict change of representation}
If \(P\subseteq D\) contains \(\bot\), is closed under suprema
of increasing \(\omega\)-chains, and \(f(P)\subseteq P\), then
\(\operatorname{fix}(f)\in P\). If \(h:D\to E\) is strict
continuous, \(f,g\) are continuous endomaps of pointed dcpo's, and
\(hf=gh\), then \(h(\operatorname{fix}(f))=\operatorname{fix}(g)\).
The assumptions do not reduce to monotonicity, or to an arbitrary
not-necessarily-strict representation map
\cite[Lemmas~2.1.20--21, p.~17]{AbramskyJungCorrectedExpandedDomainTheory}.
\end{hbtheorem}

\begin{hbexample}{def:p10-domain-recursive-factorial}{A partial-function recursive meaning}
On \([\mathbb N\rightharpoonup\mathbb N]\), define \(\Phi(p)\)
by \(\Phi(p)(0)=1\), and, for \(n>0\), define
\(\Phi(p)(n)=n\,p(n-1)\) exactly when \(p(n-1)\) is defined.
This is continuous: every produced graph pair uses at most one graph
pair from its input. Its \(k\)-th approximation defines exactly
\(0,\ldots,k-1\) with their factorial values, so its least fixed
point is total factorial. In contrast, the recursive equation
\(p=p\) has every partial function as a solution and the least
meaning is the empty function. These are editorial worked instances
of the source's partial-function domain and least-fixed-point theorem;
their construction and approximation calculation are supplied in
Annex~\ref{annex:p10.domain.examples} and
Annex~\ref{annex:p10.domain.fixpoint}.
\end{hbexample}

\begin{hbtheorem}{thm:p10-domain-failed-premises}{proved=annex:p10.domain.failures}{Failures when the fixed-point premises are removed}
On the discrete dcpo \(\mathbb N\), successor is continuous and
has no fixed point: bottom is missing. On the pointed chain
\(0<1<\cdots<\omega<\omega+1\), the monotone map sending
\(n\) to \(n+1\), \(\omega\) to \(\omega+1\), and
\(\omega+1\) to itself has an \(\omega\)-iteration supremum
\(\omega\) which is not fixed: continuity is missing.
On \(\omega+1\), successor extended by fixing \(\omega\) leaves
the finite elements invariant, but its least fixed point is not finite:
that predicate is not admissible. Finally, a non-strict constant map
from the one-point dcpo to the top of \(\{0<1\}\) intertwines
identity maps but does not preserve their least fixed points.
The monotone discontinuous example answers the actual
Exercise~2.3.9(13), rather than denying the separate existence theorem
for monotone endomaps on pointed dcpo's
\cite[pp.~15--17,32]{AbramskyJungCorrectedExpandedDomainTheory}.
\end{hbtheorem}

\subsection{Domain constructors and expanding sequences}

\begin{hbdefinition}{def:p10-domain-type-constructors}{Lifting, coalesced sum and smash product}
Let \(I\) denote the one-point dcpo. The lifting \(D_{\bot}\)
adds a \emph{new} bottom to \(D\);
write \(\operatorname{up}(d)\) for its old elements. If \(D\)
already has bottom, \(\operatorname{up}(\bot_D)\ne\bot_{D_{\bot}}\).
For pointed \(D,E\), the coalesced sum \(D\oplus E\) identifies
their bottoms and otherwise retains tagged elements \((d:1)\),
\((e:2)\), with no cross-branch order. The smash product
\(D\otimes E\) identifies all product pairs having a bottom
coordinate; other pairs keep their coordinatewise order.
\cite[Sections~3.2.3--5, pp.~44--45]
{AbramskyJungCorrectedExpandedDomainTheory}.
\end{hbdefinition}

\begin{hbtheorem}{thm:p10-domain-constructors}{proved=annex:p10.domain.constructors}{Exact constructor objects and universal properties}
These constructions are pointed dcpo's. Lifting preserves continuity
and algebraicity; its basis is new bottom together with the lifted old
basis. Coalesced sum and smash product preserve those properties for
pointed factors, with their induced tagged and paired bases.
Coalesced sum is a coproduct for \emph{strict} continuous maps, not
for all continuous maps between pointed objects. Lifting has the
natural bijection
\[
 [D\longrightarrow E]\cong
 [D_{\bot}\mathrel{\mathop{\longrightarrow}^{\bot!}}E]
                    \qquad(E\text{ pointed}),
\]
and, for pointed \(F\), smash product has the strict currying bijection
\[
 [D\otimes E\mathrel{\mathop{\longrightarrow}^{\bot!}}F]
 \cong
 [D\mathrel{\mathop{\longrightarrow}^{\bot!}}
      [E\mathrel{\mathop{\longrightarrow}^{\bot!}}F]].
\]
Thus the strict category is monoidal closed with smash product,
unit \(I_{\bot}\), and strict function space; these are not its
cartesian exponentials.
\cite[Propositions~3.2.8--9; Sections~3.2.4--5, pp.~44--45]
{AbramskyJungCorrectedExpandedDomainTheory}.
\end{hbtheorem}

\begin{hbtheorem}{thm:p10-domain-algebraic-functions}{proved=annex:p10.domain.steps}{A genuine algebraic function domain}
Let \(D\) be algebraic and \(E\) pointed, algebraic and
bounded-complete: every bounded pair in \(E\) has a join.
For compact \(a\in D,b\in E\), define
\[
 (a\mathbin{\searrow}b)(x)=
 \begin{cases}b,&a\sqsubseteq x,\\\bot_E,&a\not\sqsubseteq x.\end{cases}
\]
The compact elements of \([D\longrightarrow E]\) are exactly
finite pointwise joins of these step functions for which the join
exists; the empty join is the bottom function. This is a pointed
bounded-complete algebraic domain. Countable compact bases in \(D,E\)
give a countable compact basis of the function domain.
This is the algebraic bounded-complete case of the source's
lattice-like function-domain construction, not an assertion that
every algebraic-domain function space is algebraic
\cite[pp.~54--57, Proposition~4.1.5]
{AbramskyJungCorrectedExpandedDomainTheory}.
\end{hbtheorem}

\begin{hbdefinition}{def:p10-domain-expanding-sequence}{Embedding--projection pairs and expanding sequences}
An embedding--projection pair \(e:D\to E,p:E\to D\) consists
of continuous maps with \(pe=1_D\) and \(ep\sqsubseteq1_E\).
These equations imply that both maps are strict when their objects
are pointed. An expanding sequence has such pairs
\((e_{mn}:D_n\to D_m,p_{nm}:D_m\to D_n)\), \(n\le m\),
with identity maps at equal indices and
\(e_{kn}=e_{km}e_{mn}\) for \(n\le m\le k\).
Consequently \(p_{nk}=p_{nm}p_{mk}\).
\cite[Definition~3.1.15; Definition~3.3.6, pp.~39,48]
{AbramskyJungCorrectedExpandedDomainTheory}.
\end{hbdefinition}

\begin{hbtheorem}{thm:p10-domain-bilimit}{proved=annex:p10.domain.bilimit}{The actual chain bilimit and its two universal properties}
For an expanding sequence of pointed dcpo's set
\[
 D=\{(x_n)_n\in\prod_nD_n:p_{nm}(x_m)=x_n\ (n\le m)\}.
\]
Its order and directed suprema are coordinatewise. Let \(p_n\)
be coordinate projection and define \(e_n:D_n\to D\) by
\[
 (e_nx)_i=\begin{cases}p_{in}(x),&i\le n,\\
                         e_{in}(x),&n\le i.\end{cases}
\]
These definitions agree at \(i=n\). They form embedding--projection
pairs and satisfy \(\bigsqcup_n^{\uparrow}e_np_n=1_D\).
The \(p_n\)'s give the limit of the projection diagram; the
\(e_n\)'s give the colimit of the embedding diagram, even for
cocones consisting of arbitrary continuous maps. Its mediators are
\[
 z\longmapsto(g_n(z))_n,
       \qquad (x_n)_n\longmapsto\bigsqcup_n^{\uparrow}h_n(x_n),
\]
for a projection cone \((g_n)\) and embedding cocone \((h_n)\).
A cocone of embedding--projection pairs \((u_n,v_n)\) is
colimiting exactly when \(\bigsqcup_n^{\uparrow}u_nv_n=1\).
If each \(D_n\) is continuous with basis \(B_n\), then
\(\bigcup_ne_n(B_n)\) is a basis of \(D\); if each is algebraic,
\(K(D)=\bigcup_ne_n(K(D_n))\). Countability is preserved for a
countable sequence of countable bases.
\cite[Theorem~3.3.7; Lemma~3.3.8; Theorem~3.3.11;
Exercise~3.3.12(15), pp.~48--53]
{AbramskyJungCorrectedExpandedDomainTheory}.
\end{hbtheorem}

\subsection{Recursive type equations with actual fold and unfold}

\begin{hbdefinition}{def:p10-domain-continuous-functor}{Local continuity and continuity on embeddings}
A functor on categories of dcpo's is locally continuous when its
action on each function dcpo is Scott-continuous, hence monotone
and preserving directed suprema. For a
mixed-variance functor \(F:\mathcal D^{\mathrm{op}}\times
\mathcal D'\to\mathcal E\), this means
\[
 F(\bigsqcup^{\uparrow}A,\bigsqcup^{\uparrow}A')
   =\bigsqcup^{\uparrow}_{f\in A,\,g\in A'}F(f,g)
\]
for correctly typed directed families of maps. A functor on
embedding categories is continuous when it preserves the actual
colimit cocones of expanding sequences. This is a preservation
condition to prove, not a synonym for the desired type solution
\cite[Definitions~5.2.1,5.2.3,5.2.5, pp.~70--72]
{AbramskyJungCorrectedExpandedDomainTheory}.
\end{hbdefinition}

\begin{hbtheorem}{thm:p10-domain-mixed-embedding}{proved=annex:p10.domain.mixed}{The typed mixed-variance passage to embeddings}
Let \(F:\mathsf{DCPO}_{\bot}^{\mathrm{op}}\times
\mathsf{DCPO}_{\bot}\to\mathsf{DCPO}_{\bot}\) be locally
continuous. On embedding pairs define
\[
 \widehat F(e,e')=F(p,e'),\qquad
                   \widehat F(e,e')^*=F(e,p').
\]
These are an embedding--projection pair, define a covariant functor
on embedding categories, and preserve synchronized chain bilimits.
The diagonal functor \(G(D)=F(D,D)\),
\(G(e)=F(p,e)\), therefore preserves expanding-sequence colimits.
The function-space functor \(F(D,E)=[D\longrightarrow E]\)
satisfies this local-continuity hypothesis by the actual composition
and pointwise-supremum laws. This passage does not assert an
ordinary covariant function-space functor on all maps
\cite[Proposition~5.2.6, pp.~72--73]
{AbramskyJungCorrectedExpandedDomainTheory}.
\end{hbtheorem}

\begin{hbtheorem}{thm:p10-domain-recursive-solution}{proved=annex:p10.domain.fold}{An equation solution constructed from an expanding sequence}
Let \(G\) be a continuous functor on pointed-dcpo embedding
categories, and choose an actual embedding \(a:A\to G(A)\).
Form \(D_n=G^n(A)\), with adjacent embeddings \(G^n(a)\),
and their bilimit \(D\). With bilimit embeddings \(e_n\) and
projections \(p_n\), the maps
\[
 \operatorname{fold}=\bigsqcup_n^{\uparrow}e_{n+1}G(e_n)^*:G(D)\to D,
 \qquad
 \operatorname{unfold}=\bigsqcup_n^{\uparrow}G(e_n)p_{n+1}:D\to G(D)
\]
are actual inverse continuous maps. In particular
\(\operatorname{unfold}e_{n+1}=G(e_n)\) and
\(G(e_n)^*=p_{n+1}\operatorname{fold}\).
The one-point pointed dcpo \(I\) has an embedding into every
pointed dcpo by selecting bottom. Starting at \(A=I\) gives
the source's \(\operatorname{FIX}(G)\). A different seed is
specified rather than silently called that canonical solution.
An isomorphism alone does not prove an operational adequacy or
initial/final universal property for an unspecified language.
\cite[Lemma~5.2.2, pp.~70--71]{AbramskyJungCorrectedExpandedDomainTheory}.
\end{hbtheorem}

\begin{hbtheorem}{thm:p10-domain-recursive-examples}{proved=annex:p10.domain.recursive-examples}{Two explicit recursive semantic structures}
For a nonempty discrete dcpo \(A\), finite and infinite words
\(A^{\le\omega}\), ordered by prefix with empty word as bottom,
form an algebraic dcpo with the finite words as its compact basis.
There is an actual isomorphism
\[
 A^{\le\omega}\cong(A\times A^{\le\omega})_{\bot}:
 \quad\epsilon\mapsto\bot,\qquad aw\mapsto\operatorname{up}(a,w).
\]
Truncation at depth \(n\) and finite-prefix inclusion are
embedding--projection pairs whose supremum is identity; they
identify this object with the actual unfolding bilimit. Finite
words here are partial prefixes, not a separately tagged terminated
list value.

For the flat natural-number domain \(A=\mathbb N_{\bot}\), the
maps \(a(x)=\text{constant }x\) and \(p(f)=f(\bot_A)\)
form an embedding--projection pair
\(A\rightleftarrows[A\longrightarrow A]\).
Iterating the mixed function-space construction yields an
\(\omega\)-algebraic bilimit \(D\) with explicit inverse maps
\(D\rightleftarrows[D\longrightarrow D]\).
This solution is nontrivial because it contains the seed \(A\)
as a subdomain. Starting with the one-point seed instead gives the
one-point solution of that same equation. Application in the
nontrivial reflexive structure is
\(\operatorname{unfold}(d)(x)\), and abstraction is
\(\operatorname{fold}\) of the continuous curried denotation;
their beta/eta equations follow from the stated inverse maps.
These are editorial worked constructions using the source's
constructors, bilimit theorem and genuine-equation example,
not a claim that a solution of an equation is unique without
its specified seed or universal property
\cite[pp.~20,43,48--53,68--72]{AbramskyJungCorrectedExpandedDomainTheory}.
\end{hbtheorem}
